\section{Mathematics of the pipeline (slides 16--21)}
\label{sec:math}

This chapter is the mathematical core of the talk. Every slide equation is expanded from first principles:
definitions, intermediate steps, a worked example on \runningcomplaint{}, and why the design exists.
Non-math slides were covered earlier; here we stay inside the BioBERT triage chain.

%----------------------------------------------------------------------
\subsection{Tokenisation (slide 16, part 1)}
\label{sec:expl-token}

\begin{intuitbox}
Neural networks do not read English paragraphs. They read fixed-length lists of integer IDs.
Tokenisation is the map from raw complaint text \(X\) to that list.
\end{intuitbox}

\begin{whybox}
Without tokenisation there is no BioBERT input. This answers: how does text become numbers?
\end{whybox}

\begin{definition}[Tokenisation map]
\[
  X \xrightarrow{\tau} (t_0,t_1,\ldots,t_{M-1}), \qquad M \leq 128.
\]
Each \(t_i\) is a subword token (or a special token) with an integer ID in the vocabulary
\(\mathcal{V}\), \(|\mathcal{V}|\approx 28{,}996\).
\end{definition}

\paragraph{Why \(M\leq 128\)?}
BioBERT base was trained with this maximum length. Shorter complaints are padded with \texttt{<PAD>};
longer ones are truncated. Fixed length lets the hospital batch many requests through identical matrix shapes.

\paragraph{WordPiece algorithm.}
WordPiece keeps common clinical words whole (``fever'', ``pain'') and splits rare or misspelled forms
into pieces, e.g.\ ``headache'' \(\to\) \texttt{head} + \texttt{\#\#ache}.
Each piece occupies its own position, so \(M\) counts \emph{subword tokens}, not raw English words.
That keeps \(\mathcal{V}\) finite while still representing novel spellings at intake.

\paragraph{Three special tokens.}
\begin{center}
\small
\begin{tabular}{@{}clp{8.5cm}@{}}
\toprule
Token & ID & Purpose \\
\midrule
\texttt{[CLS]} & 101 & Classification summary; always at position \(0\) \\
\texttt{[SEP]} & 102 & End-of-sentence boundary before padding \\
\texttt{<PAD>} & 0 & Fill to length \(128\); attention mask zeros these positions \\
\bottomrule
\end{tabular}
\end{center}

\paragraph{Worked example (running complaint).}
For \runningcomplaint{} (simplified tokenisation as in the thesis):

\begin{center}
\footnotesize
\begin{tabular}{@{}cllc@{}}
\toprule
\(i\) & Raw & Token & ID \\
\midrule
0 & n/a & \texttt{[CLS]} & 101 \\
1 & I & i & 1045 \\
2 & have & have & 2031 \\
3 & a & a & 1037 \\
4--5 & headache & head + \texttt{\#\#ache} & 7994, 8772 \\
6 & and & and & 1998 \\
7 & feel & feel & 2514 \\
8--9 & feverish & fever + \texttt{\#\#ish} & 9643, 6804 \\
10 & . & . & 1012 \\
11 & n/a & \texttt{[SEP]} & 102 \\
12--127 & n/a & \texttt{<PAD>} & 0 \\
\bottomrule
\end{tabular}
\end{center}

Notice: ``headache'' and ``feverish'' each cost \emph{two} positions. That is why we say \(M\) counts subwords.

\begin{saybox}
``We split the English complaint with WordPiece, insert \texttt{[CLS]} and \texttt{[SEP]}, and pad to 128 IDs. That list is the only thing BioBERT can read.''
\end{saybox}

%----------------------------------------------------------------------
\subsection{The \texttt{[CLS]} summary vector (slide 16, part 2)}
\label{sec:expl-cls}

\begin{intuitbox}
Sequence classification needs one fixed-length vector for the whole complaint.
Averaging every token would weight ``I'' and ``and'' the same as ``feverish''.
BERT instead uses an artificial summary token \texttt{[CLS]} at position \(0\).
\end{intuitbox}

After \(L=12\) encoder layers, every token row has been updated by attending to every other token.
The row at position \(0\) is
\begin{equation}
  h_{\mathrm{[CLS]}} = H^{(L)}_{0,:} \in \mathbb{R}^{768}.
  \label{eq:hcls}
\end{equation}
The five-class head reads \emph{only} this vector. \texttt{[SEP]} marks where real text ends so padding is not treated as symptoms; \texttt{[SEP]} itself is not used for acuity prediction.

\begin{remark}
Self-attention updates the \texttt{[CLS]} row at every layer, so by layer 12 it encodes information from all symptom tokens. That is why one 768-vector can represent the full clinical wording of \(X\).
\end{remark}

\begin{saybox}
``After twelve layers we take the first row---\texttt{[CLS]}---as a 768-number summary of the whole sentence.''
\end{saybox}

%----------------------------------------------------------------------
\subsection{Embedding lookup and the three-sum (slide 17)}
\label{sec:expl-embed}

\begin{intuitbox}
A token ID is only an index into a dictionary of meaning vectors.
We also tell the model \emph{where} the token sits and which segment it belongs to.
\end{intuitbox}

\begin{whybox}
Explains how discrete IDs become continuous vectors before attention starts.
\end{whybox}

\paragraph{Word embedding.}
The matrix \(E_{\mathrm{word}} \in \mathbb{R}^{V \times 768}\) maps each ID to a dense vector:
\begin{equation}
  \mathbf{e}_i = E_{\mathrm{word}}[\mathrm{id}(t_i)] \in \mathbb{R}^{768}.
\end{equation}
After PubMed/PMC pre-training, BioBERT places clinically related terms nearer in this space than a general English model would \citep{lee2020biobert}.

\paragraph{Position and segment.}
\begin{equation}
  E_{\mathrm{pos}}(i) \in \mathbb{R}^{768}, \qquad E_{\mathrm{seg}}(s) \in \mathbb{R}^{768}, \quad s=0
\end{equation}
(for a single chief complaint, segment is always \(0\)).

\paragraph{Full input vector (element-wise sum).}
\begin{equation}
  E(t_i)_d
  = E_{\mathrm{word}}(t_i)_d + E_{\mathrm{pos}}(i)_d + E_{\mathrm{seg}}(0)_d,
  \qquad d=1,\ldots,768.
  \label{eq:embed-sum}
\end{equation}

\begin{center}
\small
\begin{tabular}{@{}lp{9.5cm}@{}}
\toprule
Symbol & Meaning \\
\midrule
\(E_{\mathrm{word}}(t_i)\) & Meaning of the token (dictionary lookup) \\
\(E_{\mathrm{pos}}(i)\) & Learned vector for position \(i\) (order matters) \\
\(E_{\mathrm{seg}}(0)\) & Segment id (here always the single-complaint segment) \\
\(E(t_i)\) & Final 768-d input for that position \\
\bottomrule
\end{tabular}
\end{center}

\paragraph{Tiny numerical sketch (first three dimensions only).}
For token ``head'' at position \(i=4\), suppose (illustrative numbers):
\begin{align*}
  E_{\mathrm{word}} &= [0.31,\,-0.18,\,0.52], \\
  E_{\mathrm{pos}}(4) &= [0.05,\,0.02,\,-0.01], \\
  E_{\mathrm{seg}}(0) &= [0.00,\,0.00,\,0.00], \\
  E(t_4) &= [0.36,\,-0.16,\,0.51].
\end{align*}
The same addition repeats for dimensions \(d=4,\ldots,768\).

\paragraph{Stacking into \(H^{(0)}\).}
Place each \(E(t_i)\) as row \(i\) of the input matrix:
\begin{equation}
  H^{(0)} \in \mathbb{R}^{M \times 768}.
  \label{eq:H0}
\end{equation}
This is the starting hidden state for encoder layer \(1\).

\paragraph{Why these dimensions?}
\begin{itemize}
  \item \(M\leq 128\): sequence length (how many tokens).
  \item \(768\): hidden width of BioBERT base (how rich each token representation is).
  \item So \(H^{(0)}\) has at most \(128\times 768 = 98{,}304\) entries per complaint---a fixed tensor shape for batched inference.
\end{itemize}

\begin{remark}
The embedding table is like a dictionary: each token ID points to 768 numbers encoding meaning.
Adding position tells the model where the word appears; adding segment (zero here) reserves the mechanism used in other BERT tasks with two sentences.
\end{remark}

\begin{saybox}
``Each token becomes 768 numbers: word meaning plus position plus segment. Stack the rows into \(H^{(0)}\) and feed the encoder.''
\end{saybox}
